\begin{table}[htbp]
\centering
\caption{Karakteristik empiris \emph{patient digital twin} dari 80 DT klinis yang ditinjau, menyoroti dominasi aliran data satu arah (angka faktual diadaptasi dari \textcite{drummond2024}).}
\label{tab:karakteristik-dt}
\small
\begin{tabular}{@{}p{6cm}p{2cm}p{5.3cm}@{}}
\toprule
\textbf{Karakteristik aliran/sinkronisasi} & \textbf{Proporsi} & \textbf{Implikasi bagi \emph{positioning}} \\
\midrule
Definisi DT yang menyebut komunikasi dua-arah & 15\% (8/55) & Mayoritas literatur pun tak mensyaratkan dua-arah \\[2pt]
DT tanpa aliran fisik$\rightarrow$virtual & 16\% (13/80) & -- \\[2pt]
DT satu-arah (fisik$\rightarrow$virtual, mirip \emph{shadow}) & 73\% (58/80) & Pola dominan; setara posisi sistem ini \\[2pt]
DT dua-arah penuh & 11\% (9/80) & Minoritas; itu pun 6/9 via rekomendasi, bukan aktuasi \\[2pt]
DT pada fase praklinis & 98\% (78/80) & Bidang masih tahap awal; \emph{proof-of-concept} wajar \\
\bottomrule
\end{tabular}
\end{table}
